% Research contribution: AX11G command architecture and relative iPRA proof.
% This is a fragment to include in the main research report.
\section{AX11G: four operands and explicit generalization}
\label{sec:ax11g}

The construction combines two changes that have not previously occurred in
one delivered table. The axiom command saves the generalization
$U(c,h)$ with its explicit first operand $c$ as binder, and it then pushes
its guarded top and continues through the complete MP parameter reader.
Consequently AX has four operands and eleven saved entries. The first
change is inherited from PA194G; the second is inherited from PA194AX11.
The contribution here is their exact combined source, a new literal table,
its command compiler, and the relative weak-arithmetic proof below.

The table has $194$ working states, $388$ transitions, dense initial state
$0$, and terminal halt label $194$. It uses a two-way infinite binary tape,
initially all zero, with the head at coordinate zero. Its dense-table hash is
\begin{quote}\small\ttfamily
7e1c6442816e5df05e961f2057719252a1976c9a195a786978bf1e2d3c8e466d.
\end{quote}
The count is $134+67-7=194$: fixed controls, raw AX readers, and simultaneous
contextual folds. It is not a sub-$194$ result.

\subsection{The exact numerical interface}

Put $p(a,b)=(a+b)(a+b+1)/2+b$ and use
\[
 I(A,B)=4p(A,B),\quad M(a,b)=4p(a,b)+1,\quad
 U(a,B)=4p(B,a)+2,\quad E(a,b)=4p(a,b)+3.
\]
Code $0$ is truth and code $1=M(0,0)$ is the contradiction target.
Write $N(B)=I(B,1)$, $A\wedge B=N(I(A,N(B)))$, and
$\exists_a B=N(U(a,N(B)))$. The Cantor stack notation $[a;w]=p(a,w)$
is total; in particular $[0;0]=0$ and both projections of zero are zero.
The separate-head numerical calculus has truth, MP, generalization, and
\begin{align*}
 h_1(D,H)&=I(\exists_D U(D,H),H),\\
 h_2(D,A,H)&=I(U(D,U(A,H)),U(A,U(D,H))),\\
 h_3(D,A,H)&=I(M(H,D),I(E(H,A),M(A,D))),\\
 h_4(D,A,H)&=I(E(H,D),I(M(A,H),M(A,D))),\\
 h_5(D,H)&=\exists_D E(H,D),\\
 h_6^Q(D,A,H)&=I(U(D,I(A,H)),I(U(D,A),U(D,H))),\\
 \rho(D,A,H,u)&=I(I(I(H,u),D),I(I(D,H),I(A,H))),\quad u\le D,\\
 \Gamma(A)&=I(A,U(A,A)).
\end{align*}
For its last head put
\begin{align*}
 P(a,b,c)&=I(c,M(c,b)),\\
 V(a,b,c)&=I(M(a,b),U(c,M(a,b))),\\
 S(a,b,c)&=U\bigl(b,I(M(a,b),P(a,b,c))\wedge I(M(c,b),c)\bigr),\\
 Q(a,b,c)&=U(c,I(S(a,b,c),N(P(a,b,c)))),\\
 g(a,b,c)&=\begin{cases}N(I(V(a,b,c),Q(a,b,c))),&c\ne a,b,\\
                         0,&\text{otherwise.}\end{cases}
\end{align*}
This is the same calculus as PA194G. The exchanged distribution head is interderivable with the primary
distribution head by the uniform formula
$(A\to(B\to C))\to(B\to(A\to C))$, used in both directions.
The delivered native certificate has $18,897$ proof lines: $8986$
instances of the restricted $\rho$ axiom and $9911$ MP steps. It assumes
neither K nor S as an extra axiom. Each numerical side condition has an
explicit component path of length at most six, proving $u\le D$ using
$x,y\le I(x,y)$. The separate flat checker verifies all these paths and
the exact exchange conclusion on a $29,932$-node formula DAG. It does not
trust the Lean-source parser used to produce the certificate. Substituting
number codes into this fixed finite proof and checking its prefix are
primitive recursive operations with quantifier-free correctness proofs;
thus the conversion is available in iPRA. No interpretation of arbitrary
arithmetic truth is involved.

A store is $(t,w,d,h)$, with theorem top $t$, stack $w$, arbitrary offset $d$,
and theorem cache $h$. Given operands $c,x,y,z$, let $D=d+x$ and $H=h+y$.
Define, in top-first order,
\[
\begin{array}{c|l}
q_0&g(H,D,c)\\q_1&\Gamma(c)\\q_2&h_6^Q(D,c,H)\\q_3&h_5(D,H)\\
q_4&h_4(D,c,H)\\q_5&h_3(D,c,H)\\q_6&h_2(D,c,H)\\
q_7&h_1(D,H)\\q_8&\rho(D,c,H,d)\\q_9&U(c,h).
\end{array}
\]
The literal command effects are
\begin{align*}
 \mathsf A_{c,x,y,z}(t,w,d,h)
   &=(0,[q_0,q_1,\ldots,q_9,t;w],z,0),\\
 \mathsf M_x(t,w,d,h)
   &=\begin{cases}(d+x,w,d+x,0),&I(t,d+x)=h,\\
                   (0,w,d+x,0),&I(t,d+x)\ne h;\end{cases}\\
 \mathsf S(t,w,d,h)
   &=\begin{cases}\mathsf{halt},&t+h=1,\\
        (h,\operatorname{tl}w,d,\operatorname{hd}w),&t+h\ne1.
       \end{cases}
\end{align*}
The fourth AX operand is real. It is consumed by the full MP reader after
AX has cleared $d,h$ and pushed $q_0$. At that point the MP input is
$(0,w',0,0)$. For $z=0$ its comparison succeeds and returns zero; for
$z>0$ it fails and returns zero. Both branches leave offset $z$.
The compiler deliberately chooses $z=0$; soundness permits every $z$.

\subsection{Total stack words and the shorter generalizer}

Let $A_0=\mathsf A_{0,0,0,0}$ and define
\[
 \mathsf{Push}=A_0\mathsf S^{10}\mathsf M_0,\qquad
 \mathsf{Pop}=A_0\mathsf S^{12},\qquad
 \mathsf{Discard}=\mathsf M_0.
\]
Then, for arbitrary natural entries,
\begin{align*}
 \mathsf{Push}(t,w,d,h)&=(0,[t;w],0,0),\\
 \mathsf{Pop}(t,[a;w],d,h)&=(t,w,0,a),\\
 \mathsf{Discard}(t,w,0,h)&=(0,w,0,0).
\end{align*}
These are total identities, including $t=1$ in Pop and all zero-stack aliases.
Indeed, $q_0,\ldots,q_8$ are divisible by four, whereas $q_9=U(c,h)\ge2$.
The shifts through the generated heads cannot pass the sum-one test.
The last test of Pop has one summand $U(0,h)\ge2$, so the saved arbitrary
top can safely be restored. Push's final comparison fails by the distinct
implication and universal residues, and clears both visible formulas.

For $j=0,\ldots,8$, apart from the additional offset needed to select a
nonzero-$u$ rho instance, the axiom selector is
\[
 \mathsf{Discard};\mathsf A_{c,D,H,0};\mathsf S^{j+2};
                         \mathsf{Pop}^{9-j}.
\]
From $(t,w,0,0)$ it returns $(q_j,w,0,0)$. Every generated selector top
has residue zero, but Pop's stronger totality means no extra hypothesis on
it is needed. For a desired rho instance with $u\le D$, use
\[
 \mathsf{Discard};\mathsf M_u;\mathsf A_{c,D-u,H,0};
                      \mathsf S^{10};\mathsf{Pop}.
\]
Offset setting satisfies $\mathsf M_i(0,w,0,0)=(0,w,i,0)$ for all $i$.

The generalization word is
\[
 \boxed{\mathsf{Gen}_i=
      \mathsf{Push};\mathsf S;\mathsf A_{i,0,0,0};\mathsf S^{11}.}
\]
From $(a,w,0,0)$, its first two parts give $(0,w,0,a)$. The axiom call
saves $U(i,a)$ as $q_9$, and eleven shifts return $(U(i,a),w,0,0)$.
There is no offset-setting command. This difference is observable:
with incoming $d=h=0$ and operand $c=1$, the saved universal is
$U(1,0)=10$, whereas the offset-binder AX11 command saves $U(0,0)=2$.

\subsection{Derivation compilation and its weak induction bound}

For an MP node with major premise $I(a,b)$ and minor premise $a$, define
\[
\begin{split}
 C(\mathrm{MP}(e,f))={}&\mathsf{Discard};\mathsf{Push};
 C(e);\mathsf{Push};C(f);\mathsf{Push};\mathsf S;\mathsf S;
 \mathsf M_b;\mathsf{Pop}.
\end{split}
\]
The initial pair of commands reserves a zero. If the recursive calls
return normally, the store before the two shifts is
$(0,[a,I(a,b),0;w],0,0)$. The first shift gives
$(0,[I(a,b),0;w],0,a)$. The second either halts with $a=1$, or gives
$(a,[0;w],0,I(a,b))$. MP then returns $(b,[0;w],b,0)$ and the total
Pop returns $(b,w,0,0)$. This includes $b=1$. Truth uses Discard,
axioms use the displayed selectors, and generalization appends
$\mathsf{Gen}_i$ to the recursively compiled premise.

\paragraph{An explicit bounded-input invariant.}
The statement ``for every stack'' must not be used as the induction
formula in iPRA. Here is a quantifier-free replacement. Fix a checked
finite proof DAG, a postorder table of its nodes, a height bound $F$ (leaves have height one, and a rule node has one
plus the largest premise height), and a requested initial top and stack $(t,w)$. Let $K\ge1$ bound every
conclusion code in the table. Put
\[
 R_0=\max(1,K,t,w),\qquad R_{j+1}=(K+R_j+3)^4.
\]
All these functions are primitive recursive. At a recursive MP call,
the only possible predecessor stacks are $[0;y]$ and $[I(a,b),0;y]$.
For $y\le R$ and $I(a,b)\le K$ their codes are at most
$(K+R+3)^4$: apply $p(v,x)\le(v+x+1)^2$ twice. Their tops are zero.
Generalization queries its predecessor at the same store.

For $k\le F$, the invariant says that every proof-table node of height
at most $k$ realizes its conclusion from every $x,y\le R_{F-k}$:
it either halts or returns the stated conclusion with stack $y$ and both
auxiliary visible entries zero. The table-node and input quantifiers are
bounded. Their finite conjunction is therefore the value of a
primitive-recursive Boolean function, and the invariant is an equation.
In the successor step $k\to k+1$, the requested outer inputs are
bounded by $R_{F-k-1}$ and the induction hypothesis permits predecessor
inputs up to $R_{F-k}$. At a leaf use its local word identity. At an MP
node, the displayed pairing bound places both recursive inputs within
that preceding-stage allowance. At a generalization node, monotonicity of $R_j$ suffices.
Guarding by $k\le F$ makes the induction equation total.
Quantifier-free induction proves the desired final instance; only then
are the originally fixed proof and store parameters universally generalized.
This proves the compiler correctness in iPRA without induction on an
unbounded all-store formula.

A coarse command-length bound is also immediate. A selector has at most
$122$ commands, a generalization suffix has $25$, and an MP node adds
at most $53$ commands to its two predecessor words. Thus
$L_0=0$, $L_{k+1}\le2L_k+128$, and $L_k\le128(2^k-1)$ is a safe common
bound. The length bound is not a bound on numerical formula codes or on
physical running time.

\subsection{The scheduler and soundness}

The nonempty continuation encoding is $e(c::v)=1+b(c,e(v))$, with
$e(\epsilon)=0$. The first command of a fresh candidate omits the outer
one. For the four-operand axiom command,
\[
 b(\mathsf A_{c,x,y,z},r)=p(c+2,p(y,p(x,p(z,r)))).
\]
The other tags are the inherited Shift and MP tags. Total Cantor
projections recover all four AX fields in the order $y,x,z$ after reading
$c+2$. An old three-operand encoding is incorrect here: it would consume
the first component of the alleged continuation as $z$.

For a derivation $e$ of code one, the word
$A_0;C(e);\mathsf S$ halts from every incoming store. The initial AX
clears $d,h$; the compiler either halts earlier or returns
$(1,w',0,0)$, and the last shift halts. Prefixing this word by
$\mathsf A_{k,0,0,0}$ gives the cofinal family of fresh candidate codes
\[
 p(k+2,p(0,p(0,p(0,e(A_0;C(e);\mathsf S))))).
\]
Each prefix terminates and leaves an arbitrary store to which the same
compiler applies. Hence the scheduler's carried store is harmless.
The candidate continuation strictly decreases; finite iteration up to a
chosen candidate and the known finite command words gives
primitive-recursive halting witnesses. No unbounded search is needed in
the proof that a supplied refutation eventually causes a halt.

Conversely, annotate $t,h$ and each live canonical stack entry with a
numerical-calculus proof of its code. Zero has its truth proof. AX adds
axiom instances, the old top, and the theorem generalization $U(c,h)$.
Its final MP reader only writes truth, irrespective of $z$. Successful
MP composes two proofs; failure writes truth. Shift moves annotations
and its halt test $t+h=1$ yields an annotated code one. Cantor zero
collapse is covered by the zero truth proof. The annotation update and
proof concatenation are primitive recursive. A bounded proof-prefix
invariant, expressed by the Boolean conjunction of proof-checker
successes, is quantifier-free; its induction therefore takes place in
iPRA. This is syntactic derivation preservation, not a truth predicate.

Consequently iPRA proves, for the exact command scheduler,
\[
 \operatorname{Halt}(R_{\mathrm{AX11G}})
 \ \longleftrightarrow\
 \exists e\,\operatorname{Prf}_{\mathcal C^Q}(e,1).
\]
The same finite, bounded-input proof applies to the displayed PA194G
command compiler: its selectors and generalizer have one fewer final
shift, its AX has three operands, and its all-store identities were
separately checked. Thus the two schedulers have iPRA-provably equivalent
halting problems. Composition with the finite propositional exchange
certificate identifies their calculus with the primary separate-head
calculus; this relative step does not require an internal translation of
arbitrary PA induction axioms.

\subsection{Literal source binding and physical scope}

Relative to PA194G, exactly the final AX macro changes from
\texttt{breakTo 1} to \texttt{node49}, the ordinary push. The independently
resolved source comparison removes old primitive jump $3553$ at path
\texttt{10101001111} and inserts the $19$-site push interval
$[3553,3572)$. Every retained primitive operation and every retained
successor agrees under relocation. Every exit of the added push enters
site $3572$, the complete MP operand reader. The old jump had skipped
that reader to old site $3641$. This is why the fourth AX field is required.

The resulting source has $3672$ primitive sites and $4779$ branch outcomes.
All $28$ simultaneous two-cell fold cases pass. The construction checker
checks every source outcome over arbitrary natural registers, arbitrary
finite nonempty spare lists and an affine gap: $1,483,363$ symbolic steps,
maximum $470$. The whole finite-support blank tape reaches the source
seed at site $25$, dispatcher $84$, head $-2$, logical registers
$(2,0,\ldots,0)$, and $26$ one-valued spare blocks after $58,154$ steps.
The checks are universal symbolic contracts and a complete finite startup,
not a bounded-run extrapolation of nonhalting.

For physical internalization, represent a finite binary tape by its two
natural-number zipper halves. Literal steps and concatenation identities
are primitive recursive. Repeated scans of $b^n$, register-list traversals,
and each fixed symbolic trace have quantifier-free induction proofs on
$n$ or on list length, with arbitrary outside contexts as parameters.
Positive local durations compose by primitive recursion. If a machine
halts at time $T$, then $T+1$ unhalted source steps with positive durations
would produce a nonhalting physical frame after that halt, contradicting
determinism and the absorbing halt convention. Bounded minimization
extracts the first source halt. Both witness transformations are primitive
recursive, and the relevant finite-prefix tests are decidable.

The separate literal-tape replay checks all $4779$ outcomes in
$1,083,079$ symbolic steps, with maximum $379$, minimum positive duration
constant three, four list lemmas derived from the actual table, and the
whole $58,154$-step startup. Its two independent source resolvers also
agree with the complete delivered source graph. These local arguments prove in iPRA the halting equivalence of the literal
AX11G table and its stated scheduler. Together with the corresponding
primary-table contracts and the foregoing source-calculus translations,
this gives the relative theorem
\[
 \mathsf{iPRA}\vdash
 \bigl(\operatorname{Halt}(M_{\mathrm{AX11G}},\mathsf{blank})
   \leftrightarrow\operatorname{Halt}(M_{194},\mathsf{blank})\bigr).
\]
The proof is ordinary computer-assisted mathematics; no new end-to-end
Lean or Rocq endpoint is asserted. A consistency calibration for the
primary table is a separate hypothesis. If that table is calibrated to
$\operatorname{Con}(\mathsf{PA})$ in iPRA, intuitionistic logic immediately
transfers the nonhalting calibration to AX11G. This relative construction
does not by itself supply the repository's missing PA-to-calculus compiler.
